% !TEX root = ../main.tex

\chapter{总结与展望}


\section{工作总结}


本文围绕跨链资产交换问题，设计并实现了一个基于公证人协调与哈希时间锁的原型系统。系统采用 Alice - ManagerA - Notary - ManagerB - Bob 的角色模型，将跨链流程抽象为用户请求、资金池锁定、公证中继、目标链响应和资产释放之间的协作关系，使各环节的控制边界清晰可分析。在机制层面，系统针对不同链的脚本表达能力提供了两种条件锁定方式：对于 EVM 类合约链，通过链上 HTLC 状态机实现 hashLock 条件领取和 timeLock 超时退款；对于 Monero、ZCash 屏蔽交易等无脚本链，由于 HTLC 的条件分支无法在链上表达，系统通过 Schnorr adaptor signature 将条件锁定放入签名补全过程，使仅具备签名能力的链也能参与原子交换。在安全性分析层面，本文针对 Notary 和 Manager 分别建立攻击模型，论证了系统在拒绝服务、参数篡改、重放攻击、抢跑攻击和密钥泄露等场景下的防御能力，并分析了两种机制在原子性保证上的安全等价性和跨机制秘密传递的正确性条件。在工程实现层面，项目完成了 Solidity 智能合约、Python 后端协调服务、scriptless 场景的 adaptor signature 模拟模块（以 BTC 作为替代实现）和前端可视化页面的组合：在本地 Hardhat 环境验证了功能正确性、安全约束、异构适配以及 gas 与存储代价，并在 Sepolia 与 Conflux eSpace 测试网上完成了真实环境的端到端跨链交换，测得各阶段确认时间、端到端耗时与 Notary 活跃时长。


\section{系统不足}


鉴于时间上的客观原因，项目在原型实现的真实性上留有遗憾。首先，无脚本链适配模块以机制模拟为主，目前以 BTC 的 secp256k1 曲线作为替代实现，尚未接入 Monero、ZCash 等真实无脚本链的交易流程，adaptor signature 的正确性仅在本地 Python 环境中验证。其次，Notary 当前更接近单节点协调模型，虽然合约和 timeLock 可以限制其直接资产风险，但其可用性仍可能影响跨链体验。再次，Manager 资金池的流动性管理较为简单，尚未充分考虑复杂汇率、手续费、滑点、限额和风险控制。此外，测试网时间代价实验仅为单次真实运行的记录，受公共 RPC 与测试网负载影响，尚缺乏多轮、多负载条件下的统计基准。最后，系统安全性主要通过攻击模型分析和实验验证说明，尚未进行形式化证明和完整安全审计。


\section{后续展望}


后续工作可以围绕安全性增强、真实环境接入和系统能力扩展三个方向推进。

在安全性增强方面，当前 Notary 为单节点协调模式，后续可以引入多 Notary 节点，通过多签触发或门限签名机制使跨链消息转发不再依赖单一节点的可用性。同时，HTLC 与 adaptor signature 的安全性可以通过形式化建模进一步论证，从而更严格地分析系统在延迟、抢跑、节点失效和参数错配等攻击场景下的安全边界。

在真实环境接入方面，无脚本链适配模块需要从本地模拟扩展到真实链的交易流程，将 adaptor signature 应用于 Monero 或 ZCash 屏蔽交易等真实无脚本链的签名与花费过程。合约链侧已在 Sepolia 与 Conflux eSpace 测试网完成单次端到端验证，后续需扩大到多轮、多负载条件下的统计测量，以给出更稳健的确认延迟、gas 消耗与 Notary 活跃时长基准，并评估不同 timeLock 参数设置的合理性。

在系统能力扩展方面，Manager 资金池可以加入流动性监控、动态费率、额度限制和异常告警机制，以提高系统面对不同交易规模和市场波动时的稳定性。随着 Solana、Cosmos 等更多异构链的接入需求出现，系统还需要进一步抽象链适配接口，使不同交易模型的链能够以统一方式接入跨链协调层。

总体而言，本文系统展示了一种将公证人协调、HTLC 哈希时间锁和 Schnorr adaptor signature 结合起来的跨链资产交换方案。该方案的核心思路是：在架构层通过统一角色流程保持跨链协作的一致性，在机制层根据目标链的脚本表达能力选择适合的条件锁定方式，在安全层通过密码学绑定和超时退出保证原子性和可退款性。借助 adaptor signature 将条件锁定移出脚本层，该方案能够把原子交换的适用范围从智能合约链扩展到无脚本链，为异构区块链之间的资产互操作提供有价值的原型参考。
